Range filters in LSM-tree storage systems are built \emph{per SSTable}:
each on-disk file carries its own filter, and the relevant input size $n$
is the number of keys in one SSTable, not the database total. Production
deployments keep SSTables small --- tens to hundreds of megabytes --- for
several reasons listed below.

\begin{enumerate}
    \item \textbf{Memory and time, compounded.} The build of the filter during compaction operations
    needs random access to keys,
    so the working set --- not just the filter --- must be held in RAM during the operation.
    Filter construction itself runs at a few million keys per second:
    \citet{grafite2024} report ${\sim}7$~M keys/s for
    Grafite at $n = 2 \times 10^8$ on a single thread.
    Both RAM and time scale linearly with $n$; compacting an $n = 2^{30}$ SSTable
    claims gigabytes of RAM and minutes of CPU on a single operation.
    This single-thread cost is unavoidable: in LevelDB, Cassandra, HBase,
    and Pebble each compaction job is single-threaded by design;
    RocksDB's optional subcompactions partition a job across threads but
    emit a separate SSTable per subcompaction, so the filter for one
    output file is still built by one thread~\citep{dong2021rocksdb}.
    \item \textbf{Cloud I/O.} Loading and saving SSTables on object
    storage backends (S3, GCS) is bandwidth-bound at the per-object
    level: a single connection to S3 sustains ${\sim}55$--$95$~MiB/s
    end-to-end~\citep{durner2023cloud}, so a multi-gigabyte file
    streamed to a remote VM takes tens of seconds.
    Splitting an object across parallel uploads scales aggregate throughput
    further but competes with the live database for network bandwidth and CPU.
    \item \textbf{Blast radius.} Corruption of one SSTable file or one
    object loses exactly the keys in that file. Smaller files bound the blast.
    \item \textbf{Sharding.} Multiple smaller SSTables can be queried in
    parallel for the on-disk scans following a filter probe, migrated
    independently between nodes, and replicated at file granularity; one
    giant file cannot.
\end{enumerate}

\paragraph{Realistic per-SSTable $n$.}
The default RocksDB \texttt{target\_file\_size\_base} is $64$~MB, with
production deployments commonly tuning to $128$~MB; entry sizes in
real workloads (key + value) are roughly $60$--$150$~bytes on
average~\citep{cao2020facebook}, putting the typical per-SSTable count
at $n \in [2^{19},\ 2^{21}]$.

Two production cases push $n$ higher. Secondary-index and counter
column families (CFs --- RocksDB's per-table partitioning, with
independent in-memory write buffers (memtables) and SSTables) store all information in the
key and use the value as a placeholder: in Facebook's social-graph
workload, the \texttt{Assoc\_count} CF has $20$-byte bigint
keys with values
under $8$~B, averaging $\sim 28$~B/entry, and similar patterns
appear in other index/counter CFs~\citep{cao2020facebook}.
KV-separation deployments
(WiscKey~\citep{wisckey2016}, RocksDB BlobDB, Pebble) store only a key
plus a $10$--$20$-byte blob handle in the SSTable, with values placed
in separate blob files. In both cases the per-SSTable count rises
to $n \in [2^{22},\ 2^{23}]$.

Range filter
papers~\citep{zhang2018surf,luo2020rosetta,vaidya2022snarf,eslami2024memento,grafite2024}
report headline numbers at $n \in [2^{24},\ 2^{28}]$, but this is a
single-filter microbenchmark convention rather than a deployment
claim: SuRF, SNARF, and Rosetta integrate per-SSTable into RocksDB;
Memento targets B-Trees (WiredTiger, MongoDB's storage engine); Grafite is presented as a
stand-alone structure. The SOSD benchmark~\citep{kipf2019sosd}
provides datasets at this same scale ($200$M keys $\approx 2^{27.6}$).

\paragraph{Why not a global filter.}
A database-wide filter is faster on the query side: range filter query
time is essentially independent of $n$, so each per-file probe is
constant-time --- but a per-file design pays one such probe per
SSTable, while a single global filter needs only one. The update
side is the opposite story: LSM compaction atomically replaces SSTables,
making per-file filter rebuild touch only the keys in that compaction,
while a global filter must either rebuild over all keys or maintain
counting structures that support deletion.
GRF~\citep{wang2024grf}, a recent global range filter for LSM-trees,
achieves single-probe queries through shape encoding; production
systems --- and our benchmarks --- still adopt the simpler per-SSTable
model.

The query-side argument is also weaker than it first looks. Within a
level (L1+), RocksDB stores SSTables as a non-overlapping
\emph{sorted run}: each file carries $[\textit{smallest},
\textit{largest}]$ key metadata in the MANIFEST (RocksDB's global level-metadata file), and a query does an
$O(\log\#\textit{files})$ binary search on this metadata to identify
only the SSTables whose key range overlaps the query~\citep{rocksdb-leveled},
then probes their per-file filters~\citep{rocksdb-bloom}.

\paragraph{Our choice of $n$.}
We benchmark at three anchor scales:
\begin{itemize}
    \item $n = 2^{20}$ --- per-SSTable production workload;
    \item $n = 2^{24}$ --- large SSTable or aggregate filter;
    \item $n = 2^{28}$ --- paper-evaluation and SOSD scale.
\end{itemize}

\paragraph{Per-key budget.}
Published range filter evaluations consistently report total per-key
budgets in the $10$--$20$~BPK range, with $14$--$16$~BPK as the typical
headline configuration (Table~\ref{tab:industry-bpk}). The LSM
point-filter convention sits at the lower end of this window: RocksDB
ships Bloom filters at $10$~BPK by default~\citep{rocksdb-bloom},
reported by~\citet{monkey2017} as the cross-system standard. We adopt
the same target range.

\begin{table}[H]
    \centering
    \small
    \setlength{\tabcolsep}{5pt}
    \begin{tabular}{@{}llcl@{}}
        \toprule
        Filter                            & BPK budget & Reported FPR                                   & Notes             \\
        \midrule
        SuRF~\citep{zhang2018surf}        & $10$--$16$ & ${\sim}10^{-3}$ at $16$~BPK\textsuperscript{*} & RocksDB-anchored \\
        Rosetta~\citep{luo2020rosetta}    & $10$--$20$ & $10^{-2}$ to $10^{-4}$                         & FPR-vs-BPK sweep  \\
        SNARF~\citep{vaidya2022snarf}     & $16$       & $6.2 \times 10^{-5}$\textsuperscript{*}        & Headline result   \\
        Memento~\citep{eslami2024memento} & $\geq 12$  & $\varepsilon$-parameterized                    & Design minimum    \\
        Grafite~\citep{grafite2024}       & ${\sim}12$ & $\leq \mathcal{L}/2^{\mathrm{BPK}-2}$          & Adversarial bound \\
        \bottomrule
    \end{tabular}
    \caption{Per-key memory budgets and corresponding false positive rates
    reported in recent range filter literature.
    FPR depends strongly on range length~$\mathcal{L}$ and workload.
    \textsuperscript{*}SuRF and SNARF FPR values at $16$~BPK are both reported in
    \citet{vaidya2022snarf} for direct comparison.}
    \label{tab:industry-bpk}
\end{table}

\paragraph{FPR targets.}\label{par:fpr-targets}
We report, for each filter, the BPK required to reach two fixed FPR
targets: $10^{-2}$ and $10^{-3}$. Both targets sit inside the operating
range of published range filter evaluations: classical
filters~\citep{zhang2018surf,luo2020rosetta,wang2024grf} anchor their
headline numbers around $10^{-2}$ to $10^{-3}$, while modern optimal
filters~\citep{vaidya2022snarf,grafite2024,eslami2024memento} sweep
deeper. They also align with the LSM point-filter convention:
Cassandra's default \texttt{bloom\_filter\_fp\_chance} is $0.01$ for
all compaction strategies except Leveled Compaction Strategy (LCS)~\citep{cassandra-bloom}.

\paragraph{Query range length $\mathcal{L}$.}
We sweep $\mathcal{L} \in \{1, 16, 128, 1024\}$. This range is
production-aligned: in Facebook's UDB workload, more than $60\%$ of
iterator operations read exactly one key, the P90 is below $100$ keys,
and scans above $10{,}000$ are very rare~\citep{cao2020facebook}; the
Yahoo! Cloud Serving Benchmark (YCSB) Workload~E caps short-range scans
at $100$. Range filter
papers~\citep{vaidya2022snarf,luo2020rosetta,grafite2024,eslami2024memento}
report headline FPR in this same window: SNARF at $\mathcal{L} = 256$,
Rosetta at $\mathcal{L} = 64$, Grafite and Memento sweeping
$\mathcal{L} \in \{1, 32, 1024\}$.